% Generated draft: requires mathematical and source review.
\documentclass[11pt]{article}
\usepackage[margin=25mm]{geometry}
\usepackage{fontspec}
\setmainfont{Latin Modern Roman}
\usepackage{amsmath,amssymb,mathtools,unicode-math}
\setmathfont{Latin Modern Math}
\usepackage{enumitem,needspace,xurl,hyperref}
\setlength{\parindent}{0pt}
\setlength{\parskip}{5pt}
\setlength{\emergencystretch}{3em}
\newcommand{\N}{\mathbb{N}}
\newcommand{\Z}{\mathbb{Z}}
\newcommand{\Q}{\mathbb{Q}}
\newcommand{\R}{\mathbb{R}}
\newcommand{\C}{\mathbb{C}}
\newcommand{\F}{\mathbb{F}}
\newcommand{\A}{\mathbb{A}}
\newcommand{\PP}{\mathbb{P}}
\newcommand{\E}{\mathbb{E}}
\newcommand{\eps}{\varepsilon}
\newcommand{\dd}{\,\mathrm{d}}
\begin{document}
\section*{UnsolvedMath ID: 30004858}
\textbf{Source problem code:} catalogue entry
\subsection{Definitions and mathematical statement}
Let \((X,0)\) be a complex surface germ. The supplied record does not specify whether the germ is reduced, normal, or has an isolated singularity, nor whether its topological type means intrinsic homeomorphism type, embedded topological type, or the oriented homeomorphism type of a link.

The source introduces a good resolution \(\pi:X_\pi\to X\): a proper bimeromorphic morphism from a smooth surface, with exceptional divisor \(E=\pi^{-1}(0)\) having simple normal crossings, which is an isomorphism away from \(E\). Write \(E=\bigcup_{v\in V(\Gamma_\pi)}E_v\), where \(\Gamma_\pi\) is the dual graph. Its vertex weights are the genus \(g(v)\geq 0\) and self-intersection \(e(v)=E_v^2<0\), as stated in the supplied excerpt.

For a generic hyperplane section \(h:(X,0)\to(\C,0)\), let \(H_\pi\) denote the strict transform of \(h^{-1}(0)\). The intersection vector suggested by the excerpt is
\[
L_\pi=(l_v)_{v\in V(\Gamma_\pi)},\qquad l_v=H_\pi\cdot E_v\in\Z_{\geq0}.
\]
Here the use of the strict transform is an explicit interpretation of the source: its initial definition refers only to the zero locus, but its subsequent explanation counts strict-transform branches. The excerpt states that when \(\pi\) factors through \(\operatorname{Bl}_0X\), the strict transform of a generic hyperplane section consists of disjoint smooth curves meeting \(E\) transversely at smooth points, and \(l_v\) counts those meeting \(E_v\). The excerpt next begins to introduce intersection numbers \(p_v\) for the strict transform of a polar curve, but ends before completing that definition.

\textbf{Question recorded in the source.} Does the topological type of \((X,0)\) impose restrictions on the possible relative positions of its polar-curve families? For which classes of surface germs does the topological type determine those relative positions?

\textbf{Unresolved formulation.} This question cannot be made into a unique precise assertion from the supplied record. The original question refers to ``these families of curves'' without supplying its antecedent. The accompanying excerpt discusses generic hyperplane sections and polar curves, whereas the curated statement refers to polar-curve families alone. Neither the families and their parameter spaces nor the equivalence relation defining their relative positions is specified. In particular, the record does not say whether relative position means intersection data on a common resolution, incidence with exceptional components, contact data, or a topological type of a configuration. It also does not specify the projection or genericity conditions defining the polar curves. Thus no particular such interpretation, class of germs, or universal quantifier over projections is asserted here. The words ``restrict'' and ``determined'' remain source questions requiring these choices before a precise theorem or conjecture can be stated.
\subsection{Short English statement}
The recorded problem asks how much the topology of a complex surface germ controls the relative positions of associated curve families, and which classes have complete determination. The supplied excerpt introduces resolution and hyperplane-section data, but leaves the intended families and the meaning of relative position unspecified. This is a statement-only draft requiring human mathematical review.
\subsection{Sources}
\begin{itemize}
\item Local UnsolvedMath registry, permanent ID \(30004858\), title ``Topological Determination of Polar Curve Positions,'' source problem code \texttt{catalogue entry}. The supplied record labels the problem \texttt{open} and its statement \texttt{corrected\_verified}; these are supplied curation labels, not independent mathematical verification.
\item Supplied report citation: \emph{Singularities (2022)}, Oberwolfach Reports, DOI \texttt{10.4171/owr/2021/47}, pp.\ 2573--2639. \url{https://doi.org/10.4171/owr/2021/47}. The registry gives proposed year \(2021\); the report title supplied in the citation uses \(2022\).
\item Attribution in the supplied excerpt: ``Polar exploration of complex surface germs,'' Anne Pichon, joint work with André Belotto da Silva, Lorenzo Fantini, and András Némethi. The separate registry field for proposer is empty.
\item Supplied literature reference: Belotto da Silva--Fantini--Némethi--Pichon, ``polar exploration.'' \url{https://arxiv.org/abs/2103.15444}.
\item Supplied literature reference: ``OWR formulation of the question.'' \url{https://doi.org/10.4171/owr/2021/47}.
\item Literature assessment dated exactly \texttt{2026-08-22}: the supplied curation labels the problem open and reports that polar-exploration work proves finite topological bounds and determines substantial combinatorial data, while general exact determination of the relative positions asked about is not established. This is supplied dated curation derived from the OWR record, not an independently verified literature assessment; its status and cited claims require independent checking.
\end{itemize}
\end{document}
